\chapter{The case of autonomous robots}
\label{ch:robots}
\doublespacing
\vspace{-2.5cm}

The previous \Cref{ch:dua} introduced the Distributed Unified Architecture, the novel framework proposed in this work. Its design and implementation have been described, from the system abstraction layer to the unit template structure. The various features and functionalities of the framework, provided by either existing tools presented in \Cref{ch:tools} or custom-made components, have been described in detail and discussed.

This work is the result of research activities spanning multiple years in the field of autonomous and distributed systems, as illustrated at the very beginning of this book. The Distributed Unified Architecture is the materialization of multiple attempts and ideas whose goal was to simplify research work involving autonomous systems and control architectures, as well as single developers or small teams in which the author has had the pleasure to work during the same period. Many of the ideas, concepts, and requirements on which DUA is based stemmed from long and fruitful discussions with colleagues and friends, usually held after a research project was completed, to assess what went well and what could be improved during the development process. This Chapter and the next tell about the fruits of these practical experiences, showing some first applications of the DUA framework and the technologies that it integrates in real-world scenarios; their aim is to complete the comprehensive description of the proposed framework, discussing some of its notable applications and the preliminary results obtained so far by applying it to real scenarios.

Although DUA is now a general-purpose framework for distributed systems, its first versions were designed with a single application in mind, which is also the subject of this Chapter: autonomous robotics. As discussed in \Cref{ch:prb,ch:tools}, autonomous robots of today are intrinsically distributed systems, from both the hardware and the software perspectives. DUA was designed to offer a unified approach to the development of such systems, providing a common ground for the integration of heterogeneous components and tools and enabling the development of complex behaviors and functionalities in a modular and scalable way.

The contents of this Chapter summarize the additional features that the DUA framework provides to the autonomous robotics domain, as well as the first results obtained by applying it and its technologies to a real-world scenario. First, some robotics software packages developed for and integrated into DUA are presented. Then, to provide a concise yet explanatory example of the capabilities of the framework, two software modules developed with DUA are briefly introduced: their design and development process and the challenges it posed were among the main motivations for the creation of the framework. Finally, some preliminary results obtained by applying the technologies and solutions on which DUA is based to real-world scenarios are illustrated.

\section{Robotic software utilities: dua-utils}
\label{sec:dua-utils}
With the system abstraction layer and the unit template finalized, a single problem is left to solve. It is of practical nature, and arises only when one starts to use the framework and, in particular, the middleware tools integrated into it for robotics software development, such as ROS 2 which is introduced in \Cref{sec:ros2}. The many features of the ROS 2 middleware compose a set of services that the robotics software developer can use to build applications and modules of various degrees of complexity, and the open-source nature of the framework has allowed for a wide community to improve it over the years, reviewing existing functionalities and adding new ones. Nonetheless, the ROS 2 middleware is a complex system, and its many features can be overwhelming for a beginner or a developer who is not familiar with all of its functionalities. Moreover, some of its APIs are not as intuitive as one would expect, requiring complex code to be written to perform simple tasks which cover instead the vast majority of their use cases.

The ROS 2 community is aware of these issues, and the development effort is improving the quality of the software every day. However, over multiple years of experience and active development of ROS 2 modules, the necessity arose to create a set of utilities that would simplify some parts of the framework, wrapping its complexity and providing a more intuitive interface to the developer. While waiting to propose the same set of changes to the community and have them integrated into the official ROS 2 distribution, a process that can notoriously happen but also take a long time, the author preferred to develop the same set of utilities as packages that could be integrated in the system configuration provided by every DUA base unit, so that they could be immediately distributed and used in the development of robotics software modules based on DUA. Note that these software modules, as shown later on, are maintained as independent projects with no specific ties to DUA, so that every published work that uses them can be easily integrated into other systems, frameworks, or architectures for the benefit of the scientific community; this is also the case of the two example modules presented later on in this Chapter.

This led to the creation of the third foundational pillar of DUA: the software collection named \texttt{dua-utils} \cite{masocco_dua-utils_nodate}. Note that, although it is a collection of utilities for ROS 2 and robotics software development, it can actually host any kind of third-party software package that can be useful in the development of distributed systems, as long as it is open-source and can be built with colcon. This can be seen as a limitation but it is only due to how it is currently integrated into DUA base units: at the very end of their setup process, the last stage clones the repository into the image filesystem and builds the packages contained in it, cleaning up build artifacts afterwards. This process has two main advantages: first, when the modules are built in this way they end up being part of the DUA base unit image, so that they can be used in the development of software modules based on DUA without the need to install them manually; second, the process is automated but performed at the end of every base unit, so if one of the utilities in \texttt{dua-utils} is updated the base unit image can be rebuilt by performing only this last step, reusing previous cached layers and saving time and bandwidth. What described so far may of course change in the future, should the necessity arise to include in \texttt{dua-utils} software packages that cannot be built with colcon or that simply concern other domains of distributed systems development. The motivating idea is that \texttt{dua-utils} is a way to customize the very base layer of the DUA framework which is provided to and intended for its users, who can actively suggest and modify its contents to adapt the framework to their most frequent necessities and usage scenarios.

To give an idea of the set of improvements that DUA provides to the ROS 2 middleware, as well as a comprehensive overview of the work that has been done so far not only with DUA but also on it, the following Sections introduce some very representative utilities that are part of the \texttt{dua-utils} collection.

\subsection{dua\_app\_management}
\label{sec:dua_app_management}
\texttt{dua\_app\_management} is a collection of software modules and source files for the management of applications and processes in the Distributed Unified Architecture. This package contains modules, both compiled and source-only, to be included in DUA-based projects. It is particularly aimed at the development of ROS 2 applications, with the aim of simplifying some common tasks and wrapping boilerplate\footnote{In the context of software development, \emph{boilerplate} code refers to long code portions that must be written each time in the same way, to perform what is usually a simple and frequent operation. Programmers usually prefer either simpler APIs or wrappers, \emph{i.e.}, equivalent code items that take less options but are simpler to write, addressing the most frequent use cases.} code, especially that which involves the lifecycle of a process and requires calling the OS to do so.

The package currently contains the \texttt{ros2\_app\_manager} C++ package, which provides two libraries and the implementation of a component container (cfr. \Cref{sec:ros2_nodes}). The first goes by the name of the package itself, and is a header-only library that implements a wrapper object for the usual main thread code of ROS 2 applications, \emph{i.e.}, processes that run single nodes or collections thereof or act as component containers. The necessity for this library comes from the fact that, if one wants to write the \texttt{main} function for a ROS 2 application going a little beyond what the examples from the official documentation show, one finds after some time that the same code is written over and over again, with only the node name and the node class changing. Operations performed by this code usually include initializing the ROS MiddleWare context, which is the entrypoint towards the underlying communication middleware used by ROS 2, configuring options for the ROS 2 node, creating a proper executor for the application, and finally adding nodes to it and spinning it. The lifecycle of such an application ends when an asynchronous event, such a signal, occurs and it must shut down, stopping the executor and destroying all such objects in reverse order to correctly release resources and terminate middleware instances. The \texttt{ros2\_app\_manager} library provides a simple API to perform all these operations:
\begin{itemize}
  \item The constructor of the \texttt{ROS2AppManager} class does all the aforementioned initializations, as well as disabling I/O streams buffering for the current process to reduce latencies, which is a common practice in ROS 2 codes.
  \item The \texttt{ROS2AppManager::run} method spins the executor, and returns only when a signal is caught, which is the usual way to handle the termination of a ROS 2 application.
  \item The \texttt{ROS2AppManager::shutdown} method stops the executor and releases all resources as previously described.
\end{itemize}
The class is implemented as a template class whose template parameter is the ROS 2 executor class to use. The library also includes the relevant system headers.

Since signals are the most common way, at least on Linux systems which are the focus of DUA, to request the termination of a ROS 2 process, the user would benefit from a finer-grained level of control over their action. To clarify the description of what follows, let us briefly introduce the concept of POSIX signals. In operating systems design and programming, and in particular for systems that follow the POSIX standard, \emph{signals} are a way to notify asynchronous events to userspace\footnote{In this context, the term \emph{userspace} denotes application programs that are run and managed by an operating system kernel.} programs, a sort of software interrupt that can be delivered to a process by the kernel either on its own behalf or on behalf of another process, which may even do so due to the user's request. Different signals are identified by numerical codes, and represent different situations and errors. For example, when we press Ctrl-C to stop a program running in a terminal, the terminal tracks that key combination and sends the \texttt{SIGINT} signal to the process it is running, which is the default signal for interrupting a process. The process can then catch the signal and perform some cleanup operations before terminating, provided that it includes a handler for that signal in its code.

In ROS 2 applications, signals are not extensively used for two reasons: first, applications running on a robot should usually run as long as the robot is on, and second, they require a signal handler to be installed in the application code and some knowledge of the OS APIs. For this reason, ROS 2 usually handles signals by itself, which is good and handy but limiting in some cases. A first example scenario is if debugging is being performed to track a bug that crashes the applications: appropriately coded signal handlers have access to information concerning what was the memory state when the crash occurred, what kind of error occurred, what was the state of the application and of buses it was accessing, and so on. A second scenario is that of a prototype robot which is being tested, a situation in which applications are usually started and stopped via remote shells opened by users. Suppose that the robot is, \emph{e.g.}, a flying drone, and that such applications allow it to fly safely: it would not be desirable if, \emph{e.g.}, the WiFi link crashes and the drone follows as well because the shells that the applications were running within close abruptly; in this case, ignoring some signals like \texttt{SIGHUP} or \texttt{SIGPIPE} would be a good idea.

The second library of this package, named \texttt{ros2\_signal\_handler}, provides a POSIX-compliant signal handler module for ROS 2 applications, which replaces the default signal handler installed by the ROS 2 middleware. The library is a header-only one, and provides a class named \texttt{SignalHandler} implemented as a singleton\footnote{In object-oriented programming, the \emph{singleton} pattern is a software design pattern that restricts the instantiation of a class to a singular instance within an entire application.}. It uses the \texttt{sigaction} API to manage handlers, allowing fine-grained control over the signals that are caught and the actions that are performed when they are caught, and can be given access to the executor running in the application to stop it when a signal is caught. The library also includes the relevant system headers. The following methods are provided:
\begin{itemize}
  \item \texttt{SignalHandler::get\_global\_signal\_handler} returns a reference to the singleton instance of the class.
  \item \texttt{SignalHandler::init} installs the signal handler, which is a static method that must be called before any other method of the class. This handler can call user-defined functions when a signal is caught, and can be configured to ignore some signals.
  \item \texttt{SignalHandler::install} to install a handler for a specific signal.
  \item \texttt{SignalHandler::ignore} to ignore a specific signal.
  \item \texttt{SignalHandler::uninstall} to restore the default handler for a specific signal.
  \item \texttt{SignalHandler::fini} to reset the \texttt{SignalHandler} instance to its initial state and restore the default signal handlers for all signals.
\end{itemize}

All these features, although useful, would be useful only when a ROS 2 application is run as a standalone process, which is not the case for all robotics applications. In reality, what happens is that nodes may also be loaded into a process at run time: such process acts as a component container and nodes are denoted as components, as explained in \Cref{sec:ros2_nodes}. In this case, the container application is usually already compiled and provided as part of the ROS 2 installation. This is where the third part of this package comes in: the \texttt{dua\_component\_container\_mt} application. It is essentially a component container based on the \texttt{MultiThreadedExecutor} executor class, just like the standard \texttt{component\_container\_mt} provided by ROS 2, but with the addition of the \texttt{ros2\_app\_manager} and \texttt{ros2\_signal\_handler} libraries. The application is intended to be used while testing ROS 2 components which may be dynamically loaded and unloaded at run time on real robot prototypes, thus accounting for the potential issues such as the one described in the example scenario above. For this reason, this container application uses the \texttt{SignalHandler} object to handle many signals and ignoring some of them, and the \texttt{ROS2AppManager} object to manage the lifecycle of the application.
